\section{Conclusions and policy implications}
Across fixed-city comparisons, the observed facility indicator expands faster than recorded activity, including when the broader endpoint frame retains source-coded-zero activity. Most of the endpoint facility scale lies in cities with positive activity at both endpoints, but city counts and facility shares describe different dimensions of network coverage. City trajectories also vary over time: relative activity growth can slow or become contraction, while an early contraction category can be followed by positive growth. Annual PPML and quarterly endpoint estimates are positive and subproportional on their respective supports. The annual log-OLS interval includes unit elasticity ($p=0.260$), and the additional quarterly follow-through contrast includes zero.

Construction support and operating evaluation should therefore be treated as linked but distinct stages. Delivery can be checked against programme conditions, while the operating period should track the observed facility stock alongside recorded activity, continuity, equipment condition and the responsibilities attached to support. Information, maintenance, continuity and payment provisions define different checks; a construction milestone is evidence at one stage, not evidence of later service.

Evaluation should distinguish coverage objectives from high-throughput objectives. A city with activity growth that lags facility expansion calls for a different inquiry from a city with activity contraction, and neither is summarized by a universal target ratio. Reporting should retain network activity totals, facility-weighted intensity, city distribution, contribution accounting and recorded-week components, with each measure tied to its denominator and horizon. This alignment makes clear what the observed network records, how activity developed and which operating duties public support assigns after construction.
